\section{生命周期、并发与可观测性}

\subsection{模型所有权与租约}

运行时作为 \texttt{llama\_model::impl} 的成员，在模型映射和张量注册完成后激活，而不是作为进程全局调度器。图执行 hook 获取 move-only 租约；它不保存裸调度器指针。析构顺序为：从全局获取注册表移除运行时、拒绝新租约、等待在途租约、停止并 join 预取线程、发出最终指标、销毁地址注册表，最后才允许模型映射销毁。这一顺序的目标是避免有效模型/上下文生命周期内的 use-after-destruction，而不是宣称对违反上游生命周期约定的调用提供保护。

\subsection{并发边界}

工作请求是有界的 \texttt{(generation, layer)} 队列。工作线程在队列锁内认领请求，然后在不持有调度器互斥锁时执行 \texttt{posix\_madvise}；相同 generation 的请求不能被两个 worker 重复认领。队列满时丢弃最旧的未认领请求并计数。这将高负载时的行为从无界积压变成可观测的降级。

\subsection{严格指标契约}

每个成功 patched benchmark 进程只允许一条 schema-1 \texttt{[SLIM-ARC-RUNTIME]} 行，记录专家发出/命中/浪费字节、回收候选/调用/跳过/失败、驻留准入/回退和压力计数。解析器拒绝重复行、缺失字段、未知字段、非十进制和溢出；baseline 必须没有该行。机制测试可验证这一契约，但不能替代真实模型上的性能实验。
